\centersection{Rezultatai ir išvados}

Įvykdžius visus iškeltus uždavinius ir realizavus \gls{risc} architektūros \glsdisp{microprocessor}{mikrovaldiklio} \glsdisp{verification}{verifikavimo} \glsdisp{framework}{karkasą}, suformuluotos žemiau pateiktos pagrindinės išvados.

\begin{enumerate}
	\item \textbf{Komercinių ir atviros prieigos įrankių analizė.} Išanalizavus analogiškas komercines \gls{eda} sistemas (\textit{Cadence}, \textit{Synopsys}, \textit{Siemens}), nustatyta, kad nors jos pasižymi aukščiausiu našumu ir DI integracijomis, jų uždarumas, \glsdisp{vendor-lockin}{priklausomybė nuo tiekėjo} bei didžiulė kaina stipriai apriboja jų naudojimą akademinėje srityje. Prieita prie išvados, kad \textit{Python} kalba ir \textit{PyUVM} biblioteka grįstas atviros prieigos karkasas yra inovatyvi, lanksti ir pilnavertė alternatyva, drastiškai sumažinanti mokymosi kreivę programinės įrangos inžinieriams, norintiems įsitraukti į aparatinės įrangos \glsdisp{verification}{verifikavimą}.
	
	\item \textbf{Infrastruktūros sprendimai ir perkeliamumas.} Išanalizavus \gls{eda} įrankių diegimo specifiką, buvo nuspręsta atsisakyti tiesioginio įrankių diegimo \glsdisp{host-operating-system}{pagrindinėse operacinėse sistemose}. Karkaso aplinkos izoliavimui pasirinkta \textit{Docker}/\textit{Podman} \glsdisp{container}{konteinerizacijos} technologija, valdoma per \textit{GNU Make} skriptus. Testavimas parodė, kad šis sprendimas yra ypač vykęs -- sistema be klaidų veikė \textit{Windows}, \textit{Linux} ir net \textit{macOS ARM M4} aplinkose.
	
	\item  \textbf{Testuojamo mikrovaldiklio programinė įranga}. Sukurta \gls{cli} komandų sistema, kuri veikia kaip dar vienas papildomas verifikavimo sluoksnis, šiomis komandomis galima testuoti tiek programinę ir aparatinę mikrovaldiklio įrangas. Šias komandas galima naudoti tiek UVM testuose ir jų imitacijose, tiek realiame \gls{fpga}, komunikuojant per terminalinę aplinką. Taipogi ši programinė įranga bus naudojama ir jau pagaminto mikrovaldiklio verifikavimo eigoje.
	
	\item \textbf{Problemos sprendimas dėl nepilno standarto.} Projektavimo fazėje susidurta su problema, jog \textit{PyUVM} biblioteka dar neturi integruotos \gls{uvm} registrų abstrakcijos lygmens (\gls{ral}) posistemės. Šiai technologinei spragai užpildyti buvo pritaikytas savas metodas -- sukurtas objektinis \textit{Python} komponentas \texttt{McuMemoryMirror}, kuris per \textit{Cocotb} virtualią sąsają atvaizduoja realią \gls{dut} atmintį ir sinchronizuojasi su \glsdisp{golden-reference}{etaloniniu modeliu}. Tai leido išlaikyti \gls{uvm} architektūros vientisumą, nenaudojant perteklinio \textit{SystemVerilog} kodo bei parodė \textit{Python} kalbos lankstumą.
	
	\item \textbf{Karkaso patikra ir kodo padengiamumas.} Atliktas integracinis testavimas patvirtino sukurtos \gls{uvm} aplinkos veikimo korektiškumą. Naudojant \gls{crv} (apribotą atsitiktinę verifikaciją) ir \textit{PyVSC} įrankius, karkasas sėkmingai sugeneravo ir į \gls{dut} išsiuntė \gls{uart} transakcijų sekas, tikrintojui (\angl{Scoreboard}) teisingai fiksuojant klaidas. Pasitelkus \textit{Cadence IMC}, užfiksuota, kad parašyti testai pasiekė daugiau nei 50\% bendrąjį \gls{soc} dizaino \glsdisp{functional-coverage}{funkcinį padengiamumą} (duomenų atmintyje pasiektas 100\%, o instrukcijų -- 98.06\% padengiamumas), kas įrodo praktinę šio \glsdisp{framework}{karkaso} naudą užtikrinant skaitmeninio dizaino kokybę.
	
	\item \textbf{Darbo eigos optimizavimas.} Automatizuota kelių polangių \textit{Tmux} terminalo aplinka iš esmės pakeitė ir pagerino testavimo eigą. Naudotojams nebereikia rankiniu būdu nustatinėti aplinkos kintamųjų ar konfigūruoti \glsdisp{simulator}{imituoklių} -- visas procesas valdomas paprastomis komandomis, todėl inžinieriai gali susitelkti išskirtinai į testų logikos rašymą.
\end{enumerate}

Šiuo metu sukurta \glsdisp{verification}{verifikavimo} karkaso 0.3 versija yra pilnai funkcionali ir yra praktiškai naudojama Kauno technologijos universiteto „KTU Proto Lab“ iniciatyvose kaip pagrindinis \glsdisp{plugin}{įskiepis} \gls{risc} mikrovaldiklio projektui. Karkaso pagalba verifikuotas \gls{rtl} dizainas ir C programinis kodas jau buvo sėkmingai \glsdisp{sythesis}{susintezuotas} bei paleistas realioje fizinėje \gls{fpga} platformoje (\textit{Basys3}), patvirtinant karkase atliktų imitacijų teisingumą. Mikrovaldiklio ir jo verifikavimo karkaso projektai buvo publikuoti \textit{GitHub} platformoje po \glsdisp{apache}{Apache 2.0 licencija}.
